\documentclass[11pt]{article}

% ---------------------------------------------------------------------------
% pcf-edge-lean : peer-review companion note
% Self-contained; compiles with pdflatex (MiKTeX). No external bibliography.
% ---------------------------------------------------------------------------
\usepackage[utf8]{inputenc}
\usepackage[margin=1in]{geometry}
\usepackage{amsmath,amssymb,amsthm}
\usepackage{array}
\usepackage{booktabs}
\usepackage{microtype}
\usepackage{xcolor}
\usepackage{listings}
\usepackage[hidelinks]{hyperref}
\hypersetup{
  pdftitle={A machine-checked Lean 4 / Mathlib formalization of the M9 Newton-edge algebraic core for polynomial continued fractions},
  pdfauthor={Papanokechi}
}

\definecolor{leankw}{rgb}{0.10,0.25,0.65}
\definecolor{leancom}{rgb}{0.30,0.45,0.30}
\definecolor{codebg}{rgb}{0.97,0.97,0.95}

\lstdefinelanguage{Lean}{
  morekeywords={theorem,lemma,axiom,def,noncomputable,by,have,calc,rw,exact,
    open,namespace,end,Type,CommRing,IsDomain,Field,CharZero,Polynomial},
  sensitive=true,
  morecomment=[l]{--},
  morecomment=[n]{/-}{-/},
  morestring=[b]",
}
\lstset{
  language=Lean,
  backgroundcolor=\color{codebg},
  basicstyle=\ttfamily\footnotesize,
  keywordstyle=\color{leankw}\bfseries,
  commentstyle=\color{leancom}\itshape,
  breaklines=true,
  breakatwhitespace=false,
  postbreak=\mbox{\textcolor{leankw}{$\hookrightarrow$}\space},
  columns=fullflexible,
  frame=single,
  framesep=4pt,
  xleftmargin=4pt,
  showstringspaces=false,
  literate=
    {≠}{{${}\neq{}$}}1
    {≤}{{${}\leq{}$}}1
    {≥}{{${}\geq{}$}}1
    {→}{{${}\to{}$}}1
    {ℝ}{{$\mathbb{R}$}}1
    {ℕ}{{$\mathbb{N}$}}1
    {ℤ}{{$\mathbb{Z}$}}1
    {α}{{$\alpha$}}1
    {β}{{$\beta$}}1
    {γ}{{$\gamma$}}1
    {θ}{{$\theta$}}1
    {ξ}{{$\xi$}}1
    {Σ}{{$\Sigma$}}1
    {∘}{{$\circ$}}1
    {·}{{$\cdot$}}1
    {−}{{$-$}}1
    {₀}{{$_0$}}1
    {ₖ}{{$_k$}}1
    {ₙ}{{$_n$}}1
    {ᵏ}{{$^k$}}1
    {⁻¹}{{$^{-1}$}}1
}

\newtheorem{theorem}{Theorem}
\newtheorem{proposition}[theorem]{Proposition}
\newtheorem{definition}[theorem]{Definition}
\newtheorem{remark}[theorem]{Remark}

\newcommand{\xizero}{\xi_{0}}
\newcommand{\cone}{\{\texttt{propext},\,\texttt{Classical.choice},\,\texttt{Quot.sound}\}}
\newcommand{\Hedge}{\textsf{H\textsubscript{edge}}}

\title{\textbf{A machine-checked Lean 4 / Mathlib formalization of the\\
M9 Newton-edge algebraic core for polynomial continued fractions}\\[4pt]
\large (conditional on the analytic edge\,$\Rightarrow$\,Borel-radius step)}

\author{Papanokechi\\
\small ORCID: \href{https://orcid.org/0009-0000-6192-8273}{0009-0000-6192-8273}}

\date{June 5, 2026 \quad(\texttt{pcf-edge-lean} v0.1.0, draft for review)}

\begin{document}
\sloppy
\setlength{\emergencystretch}{3em}
\maketitle

\begin{abstract}
\noindent
We formalize, in Lean~4 / Mathlib, the \emph{purely algebraic, finitary core}
of the M9 Newton-edge lemma for polynomial continued fractions (PCFs) with unit
numerator and partial denominator $b(n)=\sum_{k=0}^{d}\beta_k n^k$
($d\ge 2$, $\beta_d\neq 0$). Four statements are proved with the clean Lean
axiom cone $\cone$ and zero \texttt{sorry}: (T1) shift-invariance of the leading
coefficient, $[\theta^d]\,b(\theta+1)=\beta_d$; (T2) the closed form of the
slope-$1/d$ edge polynomial
$\chi_d(c)=1+(-1)^{d+1}(\beta_d/d^{d})c^{d}$; (T3a) the
affine-reindexing coefficient law --- the degree-$d$ coefficient of
$b\circ(\alpha X+\gamma)$ equals $\alpha^{d}\beta_d$, which is
$\mathrm{lead}(b\circ\varphi)$ since $\deg=d$ for $\alpha\neq0$, and is
$\gamma$-independent; and (T3b) the covariance
$\xizero(\alpha^{d}\beta_d)=\alpha^{-1}\xizero(\beta_d)$ of the
Borel-singularity radius $\xizero=d/\beta_d^{1/d}$ (the T3b covariance is
formalized for $\beta_d>0$; the complex root-modulus identity $|c|=\xizero$ is
\emph{VERIFIED} in computer algebra, not Lean-formalized). The analytic step
\emph{Newton edge $\Rightarrow$ Borel singularity radius} (Wasow~\S19 sectorial
existence; Birkhoff--Trjitzinsky Borel summability) is \textbf{out of scope}:
it is encoded as an explicit, unproved axiom \Hedge{} and is \emph{not}
discharged. Consequently the Borel-radius end statement is proved
\emph{conditional on} \Hedge, with that hypothesis visible in its axiom cone by
design. This promotes the edge lemma from the source grade
``symbolically derived'' (a hand induction corroborated by computer algebra) to
machine-checked \emph{proven} for its finitary core, while leaving the analytic
input honestly labeled. The toolchain is pinned (Lean
\texttt{v4.29.0}; Mathlib rev \texttt{8a17838\ldots}); all axiom cones are
reproducible by \texttt{\#print axioms}. This note is a companion to the M9
deposit (concept DOI \texttt{10.5281/zenodo.20542161}).
\end{abstract}

\paragraph{Epistemic-status legend (used throughout).}
\textbf{PROVEN} $=$ the Lean \texttt{\#print axioms} cone is contained in
$\cone$ \emph{and} contains no \texttt{sorryAx}, with zero
\texttt{sorry}/\texttt{admit}/\texttt{native\_decide} in source.
\textbf{PROVEN CONDITIONAL on \Hedge} $=$ otherwise clean, but the cone
additionally contains the explicit axioms \texttt{edge\_implies\_borel\_radius}
and \texttt{borelRadius}. \textbf{VERIFIED} $=$ checked numerically/symbolically
(computer algebra) but not in Lean. \textbf{CITED-CONDITIONAL} $=$ relied upon
from the literature and not discharged here. These four grades are kept
distinct and the source comments match the realized axiom cones.

\section{Scope and the conditional-core discipline}

For a PCF with unit numerator and partial denominator
$b(n)=\sum_{k=0}^{d}\beta_k n^{k}$, the convergent denominators $Q_n$ obey the
Wallis recurrence $Q_n=b(n)Q_{n-1}+Q_{n-2}$. The generating function
$F(z)=\sum_n Q_n z^{n}$ is annihilated by
\begin{equation}\label{eq:operator}
  P \;=\; 1 \;-\; z\,b(\theta+1)\;-\;z^{2},
  \qquad \theta=z\,\frac{d}{dz}\ \ (\theta\,z^{n}=n\,z^{n}).
\end{equation}
At $z=0$ the Newton polygon of $P$ has a single dominant edge of slope $1/d$;
the WKB/edge balance on that edge produces the edge polynomial $\chi_d$, whose
root modulus is the Borel-singularity radius $\xizero$. The passage
\emph{edge $\Rightarrow$ analytic Borel radius} is classical analysis
(Wasow~\S19; Birkhoff--Trjitzinsky). 

This note deliberately separates two things that are easy to conflate:
\begin{itemize}\setlength\itemsep{2pt}
  \item the \emph{finitary algebra} underneath the edge --- shift-invariance,
        the closed form of $\chi_d$, and the affine leading-coefficient law ---
        which is purely a statement about polynomials and real powers, and is
        formalized here as PROVEN;
  \item the \emph{analytic} edge\,$\Rightarrow$\,radius implication, which is
        \emph{not} formalized and is taken as the explicit hypothesis \Hedge.
\end{itemize}
This is the ``conditional-core'' pattern: formalize the finitary core to a clean
axiom cone, and take the hard analytic fact as a clearly-labeled hypothesis.
The end statement is then honestly \emph{``$\xizero$ covariance, conditional on
\Hedge''}, never silently ``$\xizero$ verified''.

\section{The four proven statements}

All four are in \texttt{PcfEdge/Basic.lean}. We give the exact Lean statements
(verbatim) followed by the mathematical content.

\subsection*{T1 --- shift-invariance of the leading coefficient (gate G1)}
\begin{lstlisting}
theorem shift_leading_coeff {R : Type*} [CommRing R] [IsDomain R] (p : R[X]) :
    (p.comp (X + C 1)).coeff p.natDegree = p.leadingCoeff
\end{lstlisting}
Composition with $X+1$ preserves the top coefficient; equivalently
$[\theta^{d}]\,b(\theta+1)=\beta_d$. (The domain hypothesis is needed because
the underlying Mathlib lemmas \texttt{natDegree\_comp} / \texttt{leadingCoeff\_comp}
require \texttt{NoZeroDivisors}.)

\subsection*{T2 --- edge polynomial closed form (gate G2)}
\begin{lstlisting}
theorem edge_poly_closed_form {R : Type*} [Field R] [CharZero R] (d : ℕ) (_hd : d ≠ 0)
    (βd c : R) :
    1 - βd * (-(c / (d : R))) ^ d
      = 1 + (-1) ^ (d + 1) * (βd / (d : R) ^ d) * c ^ d
\end{lstlisting}
With $a_{00}=1$ and edge coefficient $a_{1,d}=-\beta_d$, the edge balance
$a_{00}+a_{1,d}\,(z\,u^{d})=0$ gives
\[
  \chi_d(c)\;=\;1-\beta_d\bigl(-\tfrac{c}{d}\bigr)^{d}
  \;=\;1+(-1)^{d+1}\frac{\beta_d}{d^{d}}\,c^{d},
\]
proved $d$-free: the lemma is stated for general $d$ with the single hypothesis
$d\neq 0$ (equivalently $d\ge 1$ over $\mathbb{N}$). The manipulation
$\bigl(-(c/d)\bigr)^{d}=(-1)^{d}c^{d}/d^{d}$ is valid for every $d\ge1$ in
characteristic~$0$.

\subsection*{T3a --- affine-reindexing leading-coefficient law (gate G4)}
\begin{lstlisting}
theorem reindex_leading_coeff {R : Type*} [Field R] (p : R[X]) {d : ℕ}
    (hd : p.natDegree = d) (α γ : R) (hα : α ≠ 0) :
    (p.comp (C α * X + C γ)).coeff d = α ^ d * p.leadingCoeff
\end{lstlisting}
The Lean conclusion is literally about the \emph{degree-$d$ coefficient}
$[\,\cdot\,]_d$, written \texttt{.coeff d}, \emph{not} \texttt{.leadingCoeff}.
The two coincide here precisely because $\deg(b\circ\varphi)=d$ whenever
$\alpha\neq0$: the proof first establishes
$\texttt{natDegree}(b\circ\varphi)=d$ (via \texttt{natDegree\_comp} and
$\deg(\alpha X+\gamma)=1$), so the degree-$d$ coefficient \emph{is} the leading
coefficient of $b\circ\varphi$. Hence
$\mathrm{lead}(b\circ\varphi)=\alpha^{d}\beta_d$, \emph{independent of
$\gamma$}.

\subsection*{T3b --- the $\xizero$ covariance (gate G4)}
\begin{lstlisting}
noncomputable def xi0val (d : ℕ) (βd : ℝ) : ℝ := (d : ℝ) / βd ^ ((1 : ℝ) / (d : ℝ))

theorem xi0_covariance (d : ℕ) (hd : d ≠ 0) (α βd : ℝ) (hα : 0 < α) (hβ : 0 < βd) :
    xi0val d (α ^ d * βd) = α⁻¹ * xi0val d βd
\end{lstlisting}
With $\xizero(\beta_d)=d/\beta_d^{1/d}$ (real \texttt{rpow}), the rescaling
$\beta_d\mapsto\alpha^{d}\beta_d$ scales $\xizero$ by $\alpha^{-1}$.
\textbf{Scope (corrected).} This is an identity about the \emph{real positive
value} $\xizero$ as a function of the leading coefficient; it asserts
\emph{nothing} about the roots of $\chi_d$. The root-modulus fact G3 --- that
every complex root $c$ of $\chi_d$ has $|c|=\xizero$ --- is \emph{VERIFIED} in
computer algebra (\S5) but is \emph{not} formalized in Lean (\S6). The numerical
value of $\xizero$ does equal that common root modulus, but
\texttt{xi0\_covariance} captures only the covariance of $\xizero$, not the
root-modulus statement.

\textbf{Sign regime ($\beta_d>0$).} T3b carries the hypotheses
$\alpha>0$ and $\beta_d>0$: real \texttt{rpow} is well-behaved (and the
\texttt{Real.mul\_rpow}/\texttt{Real.rpow\_mul} step valid) only for nonnegative
bases. The abstract's object-level hypothesis on a PCF is $\beta_d\neq0$, which
is strictly weaker. \textbf{The formalized covariance therefore covers
$\beta_d>0$ only}; the case $\beta_d<0$ is outside the Lean scope. This is not a
gap in the mathematics but in what is formalized: for $\beta_d<0$ the relevant
object is the common \emph{complex} root modulus
$|c|=d/|\beta_d|^{1/d}$ (real \texttt{rpow} of a negative base is defined in
Mathlib but is not the analytic continuation one wants here), which is exactly
the complex root-modulus fact G3 --- VERIFIED in computer algebra (for
$\beta_d<0$ the modulus $|c|=d/|\beta_d|^{1/d}$ was re-checked here for
$d=2..7$, all $d$ roots to $30$ digits), not Lean-formalized. A referee asking
``what happens for negative leading coefficient?'' is answered: it is handled by
the VERIFIED complex-modulus G3, deliberately not by the PROVEN
real-\texttt{rpow} T3b. This $\beta_d>0$ regime is \emph{inherited by the
conditional end statement}: \texttt{borelRadius\_covariance} carries
\texttt{hlc : 0 < p.leadingCoeff} (i.e.\ $\beta_d>0$), which is exactly the
hypothesis that both its second \texttt{edge\_implies\_borel\_radius} application
(on $b\circ\varphi$, whose leading coefficient $\alpha^{d}\beta_d$ must be
positive) and its \texttt{xi0\_covariance} call require. So for a reader who
cares only about the end statement, $\beta_d>0$ bites precisely there, through
\texttt{hlc}; it is load-bearing, not decorative.

\section{The \Hedge{} axiom and the conditional end statement}

The analytic input is isolated as two explicit declarations, so that anything
depending on it is flagged in its axiom cone:
\begin{lstlisting}
-- ξ₀ read off a real polynomial: the same real function as in T3b.
noncomputable def xi0 (p : ℝ[X]) : ℝ := xi0val p.natDegree p.leadingCoeff

axiom borelRadius : ℝ[X] → ℝ                    -- OPAQUE; not constructed

axiom edge_implies_borel_radius (p : ℝ[X]) (hd : 2 ≤ p.natDegree)
    (hlc : 0 < p.leadingCoeff) : borelRadius p = xi0 p   -- this is H_edge
\end{lstlisting}
Here $\texttt{xi0 } p$ is \emph{defined} as
$\texttt{xi0val}(\,p.\texttt{natDegree},\,p.\texttt{leadingCoeff}\,)$, i.e.\ the
same real function $\xizero=d/\beta_d^{1/d}$ from T3b, now evaluated at the
degree and leading coefficient of $p$; this is the only place $\xizero$ enters
the polynomial-level statement. $\texttt{borelRadius}$ is the analytic
Borel-singularity radius, left \emph{opaque} (we do not construct it).
$\texttt{edge\_implies\_borel\_radius}$ is exactly \Hedge: for a degree-$\ge2$
denominator with positive leading coefficient, the analytic radius equals the
algebraically computed $\xizero$. It is \emph{not} proved. The conditional
theorem is then
\begin{lstlisting}
theorem borelRadius_covariance (p : ℝ[X]) (α γ : ℝ) (hd : 2 ≤ p.natDegree)
    (hlc : 0 < p.leadingCoeff) (hα : 0 < α) :
    borelRadius (p.comp (C α * X + C γ)) = α⁻¹ * borelRadius p
\end{lstlisting}
i.e.\ $\xizero(b\circ\varphi)=\alpha^{-1}\xizero(b)$ for $\varphi(n)=\alpha n+\gamma$,
$\alpha>0$ --- \textbf{proven conditional on \Hedge}. Its proof composes
\Hedge{} (applied to $b$ and to $b\circ\varphi$) with the purely algebraic T3a
and T3b; no analytic content is smuggled in.

\section{Soundness: axiom cones, build, toolchain}

The decisive check, under the project's discipline, is the axiom cone, not a
green build (a \texttt{sorry} compiles). The cones below were obtained from
\texttt{\#print axioms} by compiling against the pinned toolchain.

\begin{center}
\small
\renewcommand{\arraystretch}{1.3}
\begin{tabular}{@{}l>{\raggedright\arraybackslash}p{6.6cm}c@{}}
\toprule
\textbf{Declaration} & \textbf{Axiom cone} & \textbf{Status}\\
\midrule
\texttt{shift\_leading\_coeff}    & $\cone$ & PROVEN\\
\texttt{edge\_poly\_closed\_form}  & $\cone$ & PROVEN\\
\texttt{reindex\_leading\_coeff}   & $\cone$ & PROVEN\\
\texttt{xi0\_covariance}           & $\cone$ & PROVEN\\
\texttt{borelRadius\_covariance}   & $\cone$ $\cup\,\{\texttt{borelRadius},$ $\texttt{edge\_implies\_borel\_radius}\}$ & COND.\\
\bottomrule
\end{tabular}
\end{center}

\noindent
The first four cones are \emph{exactly} the clean set, with no
\texttt{sorryAx}. The end statement's cone additionally exposes the two
\Hedge{} axioms --- expected and correct, and the precise sense in which the
analytic claim remains undischarged.

\paragraph{Build / hygiene.}
\texttt{lake build} completes with exit code~0. The sources contain zero
\texttt{sorry}, \texttt{sorryAx}, \texttt{admit}, or \texttt{native\_decide}
(the only textual occurrence of ``sorryAx'' is inside a documentation comment).

\paragraph{Toolchain (pinned, reproducible).}
Lean \texttt{leanprover/lean4:v4.29.0}; Mathlib rev
\texttt{8a178386ffc0f5fef0}\allowbreak\texttt{b77738bb5449d50efeea95} (inputRev \texttt{v4.29.0}),
committed via \texttt{lake-manifest.json}. The Mathlib lemmas relied upon
(each \texttt{\#check}-confirmed in the installed source) include
\texttt{natDegree\_comp}, \texttt{natDegree\_X\_add\_C}, \texttt{leadingCoeff\_comp},
\texttt{monic\_X\_add\_C}, \texttt{natDegree\_linear}, \texttt{leadingCoeff\_linear},
\texttt{Real.rpow\_natCast}, \texttt{Real.mul\_rpow}, \texttt{Real.rpow\_mul},
and \texttt{Real.rpow\_one}; no lemma name was used from memory.

\section{Numeric/symbolic ground truth (VERIFIED, gate G1--G4)}

Before formalizing, the load-bearing identities were checked in computer
algebra; all four gates were \emph{re-run in this session} (sympy for the
identities; \texttt{mpmath} for the G3 root moduli) for $d=2..6$, with $d=7$
spot-checks, so the Lean \emph{statements} match a verified gate:
\begin{center}
\renewcommand{\arraystretch}{1.2}
\begin{tabular}{@{}clc@{}}
\toprule
\textbf{Gate} & \textbf{Statement} & \textbf{Checked}\\
\midrule
G1 & $[\theta^{d}]\,b(\theta+1)=\beta_d$ & $d=2..6$\\
G2 & $1-\beta_d(-(c/d))^{d}=1+(-1)^{d+1}(\beta_d/d^{d})c^{d}$ & $d=2..6$\\
G3 & every root $c$ of $\chi_d$ has $|c|=d/\beta_d^{1/d}=\xizero$ & $d=2..6$\\
G4 & $\mathrm{lead}(b\circ\varphi)=\alpha^{d}\beta_d$;\ \ $\xizero(b\circ\varphi)=\alpha^{-1}\xizero(b)$ & $d=2..6$\\
\bottomrule
\end{tabular}
\end{center}
G3 at $d=3$ in particular settles, numerically, the previously-deferred direct
cubic check (re-confirmed here for $d=2..6$, $d=7$ spot-checked, by computing
all $d$ complex roots and verifying $|c|=\xizero$ to $30$ digits). \textbf{G3 is
deliberately \emph{not} reflected on the Lean side}: T3b is the covariance of the
real value $\xizero$, not a root-modulus statement (\S6). All four gates are
\emph{VERIFIED}, a grade strictly weaker than the PROVEN Lean core and recorded
as such.

\section{What is \emph{not} proved (scope and ``could not confirm'')}

\begin{itemize}\setlength\itemsep{3pt}
  \item \textbf{The analytic \Hedge.} The edge\,$\Rightarrow$\,Borel-radius
        implication (Wasow~\S19 sectorial existence; Birkhoff--Trjitzinsky
        Borel summability) is \emph{cited-conditional}: encoded as an axiom,
        not formalized. No analysis/asymptotics content of Mathlib was used to
        discharge it. This is the single load-bearing analytic assumption.
  \item \textbf{Full complex root-set of $\chi_d$ (G3).} We do \emph{not}
        formalize ``every complex root of $\chi_d$ has modulus $\xizero$'' in
        Lean at all --- it is \emph{VERIFIED} only in computer algebra. T3b
        (\texttt{xi0\_covariance}) formalizes the covariance of the \emph{real
        value} $\xizero$, which is a strictly different statement and must not be
        read as capturing G3. This is a deliberate scope choice, not a hidden
        claim.
  \item \textbf{No $d\ge 7$ or all-$d$ analytic claim.} The VERIFIED gate is
        $d\le 6$; the PROVEN T1--T3 are $d$-free as algebraic identities, but
        the analytic $\xizero$ interpretation rests on \Hedge{} regardless of
        $d$.
  \item \textbf{Could-not-confirm list:} empty for Mathlib lemma names (every
        lemma used was located and \texttt{\#check}ed in the installed source).
\end{itemize}

\paragraph{Honest one-line claim.}
\emph{The finitary algebra of the M9 Newton edge --- shift-invariance, the
closed form of $\chi_d$, the $\gamma$-free affine leading-coefficient law, and
the $\alpha^{-1}$ covariance of $\xizero=d/\beta_d^{1/d}$ --- is machine-checked
PROVEN in Lean~4 with a clean axiom cone; that the analytic Borel radius equals
$\xizero$ and inherits the covariance is PROVEN only conditional on the explicit,
unproved hypothesis \Hedge.}

\section{Reproducibility}

The axiom-cone table in \S4 is the load-bearing soundness claim, and --- as a
\texttt{sorry} compiles --- it is established only by an actual
\texttt{\#print axioms} run, not by a green build. To reproduce it exactly:
\begin{enumerate}\setlength\itemsep{2pt}
  \item Clone the repository and use the pinned \texttt{lean-toolchain}
        (\texttt{v4.29.0}) and \texttt{lake-manifest.json}.
  \item \texttt{lake exe cache get} (prebuilt Mathlib oleans), then
        \texttt{lake build} (exit code~0).
  \item Run the committed checker against the pinned toolchain:
        \texttt{lake env lean cone\_check.lean}. It imports \texttt{PcfEdge} and
        prints \texttt{\#print axioms} for all five declarations. The output
        below is \textbf{transcribed} from a prior-session compile and is
        \textbf{not} a live capture from the present preparation session; it is
        the value the operator's re-run is expected to reproduce:
\begin{lstlisting}
-- EXPECTED (transcribed from a prior compile; confirm via the run above).
-- NOT a live capture from the preparation session.
'shift_leading_coeff'    depends on axioms: [propext, Classical.choice, Quot.sound]
'edge_poly_closed_form'  depends on axioms: [propext, Classical.choice, Quot.sound]
'reindex_leading_coeff'  depends on axioms: [propext, Classical.choice, Quot.sound]
'xi0_covariance'         depends on axioms: [propext, Classical.choice, Quot.sound]
'borelRadius_covariance' depends on axioms:
  [propext, Classical.choice, borelRadius, edge_implies_borel_radius, Quot.sound]
\end{lstlisting}
        If reproduced, this is exactly the \S4 table: the first four cones are the
        clean set, and only the conditional end statement exposes the two
        \Hedge{} axioms.
\end{enumerate}
\textbf{Provenance of the \S4 cones (read this with the table).} The cones were
obtained by compiling \texttt{PcfEdge/Basic.lean} (\textsc{sha-256}
\texttt{cac02d77\ldots1c8dd48a7}) against Mathlib rev \texttt{8a17838\ldots} in a
\emph{prior} working session. The present session did \emph{not} re-run the
check --- Lean-toolchain execution was unavailable in its sandbox --- but it
\emph{did} re-verify, this session, that the source SHA-256 is unchanged, so the
transcribed table still refers to the bytes on disk. Until a fresh
\texttt{lake env lean cone\_check.lean} is run on a clean clone, the \S4 table
should be read as \emph{VERIFIED in a prior run and pending operator
re-confirmation}, not as freshly captured. That re-run is the \textbf{required}
final gate before tagging a release or minting a DOI.

\paragraph{Relation to the deposited record.}
This artifact is a companion to the M9 preprint (concept DOI
\texttt{10.5281/zenodo.20542161}), whose own statement notes that ``there is no
Lean formalization here'' and that the general-$k$ step is ``a hand induction
($k=0..8$ machine-checked in sympy)''. The present development closes precisely
that gap for the finitary core, while preserving the paper's
cited-conditional treatment of the analytic edge\,$\Rightarrow$\,radius step.

\paragraph{AI-assistance disclosure.}
Prepared with AI assistance (GitHub Copilot; Anthropic Claude) for Lean
drafting, computer-algebra verification, and manuscript preparation. All axiom
cones were obtained by compiling against the pinned toolchain; the author takes
responsibility for the content and for the placement of each claim within the
four-grade epistemic partition.

\end{document}
